\section{Limitations}
\label{sec:limitations}

\subsection{Identification scope}
\label{sec:lim_identification}

Identification is descriptive throughout. The two design-based
strategies that would extract a causal estimate return null at
the available bandwidths. A sharp regression discontinuity at
the convite/preg\~ao statutory caps is gated by the absence of
bunching in the McCrary density test, and a
difference-in-differences exploiting the $2018$ cap raise is
gated by the failure of parallel trends. The within-PBU coefficient
$\valPctSixfour$ on the broad sample reverses to
$\valMatchOverlapCoef$ under overlap-restricted matching and
$\valMatchPSCoef$ under propensity-score trimming; the
segment-level decomposition in \S\ref{sec:results_overlap}
attributes the aggregate reversal to the population-weighted
average of a positive Q4-tender-value coefficient and three
negative Q1--Q3 coefficients, all robust to design choice. The
sensitivity bounds (\citet{cinelli2020making} robustness value
$\valRVqOne$, \citet{oster2019unobservable}
$\hat\delta = \valOsterDelta$) discipline the broad-sample
association against selection on observables already absorbed by
item, year, and procuring-unit fixed effects, but do not extend
to selection on unobservables that the overlap restriction
surfaces. The $\valHeadlineRange$ headline range is a conditional
descriptive association, not an estimate of cover bidding's
causal effect on prices, and the
$\valPBUgradientRatio\!\!\times$ extreme-quartile gradient is
observationally consistent with the framework's detection-cost
comparative static without identifying the causal channel. Under
the screening framing of \S\ref{sec:emp_estimand}, this
descriptive scope is sufficient: the screening object is the
loser-side participation footprint where the cartel deploys it,
not a counterfactual treatment effect on prices.

\subsection{Construct scope}
\label{sec:lim_construct}

The construct flags loser-side participation patterns, not cartel
members. The validation discipline established in
\S\ref{sec:cade}---cobidders inside the always-loser stratum as
the structurally appropriate object given the data layer,
direct-defendant AUC of $\valAUCdirectStd$ as the design's
empirical signature---defines the limit. An enforcement program
that requires firm-level membership identification will need a
different statistic. Within the always-loser stratum, the
participation footprint is observationally compatible with three
non-competitive readings: cover bidding under cartel rotation,
capacity-transfer signaling for subcontracting or affiliate
consortia, and market-intelligence accumulation. The framework of
Online Appendix~A is built around the cover-bidding reading,
which is the reading the strongest heterogeneity dimensions
match (modal asymmetry in \S\ref{sec:mech_modal}, detection-regime
gradient in \S\ref{sec:results_oversight}); the bid-level bridge
in \S\ref{sec:cade_bridge} narrows it further to the
credible-cover-bidding variant. The other two readings are not
refuted, and the empirical work is organized around the
cover-bidding reading without claiming it as unique. Two finer
mechanism predictions---a cell-heterogeneity grid and a discrete
first-time-FL imprint---move in the predicted direction but the
data do not adjudicate them under disciplined estimation
(Online Appendix~C); neither is load-bearing for the screening
interpretation. Bid-aggressiveness against the reference price
remains an untested prediction of the cover-bidder type:
\textit{Valor Unit\'ario Refer\^encia} is populated for only
$\valTBBLRefCovConvite$ of convite and $\valTBBLRefCovPregao$ of
preg\~ao losing bids, leaving most firms without per-firm
coverage; that part of the bridge is reserved for jurisdictions
with universal reference-price archival.

The construct's behavior under deliberate adversarial adaptation
is bounded but not invariant
(Online Appendix~G,
Table~\ref{tab:adversarial_adaptation}). The screening statistic
is resilient to periodic rotation of cover bidders and to
occasional wins evading the bright-line filter (AUC essentially
unchanged), degrades modestly under CNPJ splitting
($\valAdaptCNPJDeltaPP$ points), and is substantially vulnerable
to threshold-aware capping ($\valAdaptCapDeltaPP$ points);
combining capping with rotation drives AUC from $\valAUCBase$ to
$\valAUCAdaptCombined$. The threshold rule is the most exposed
margin, and any application of the construct requires periodic
recalibration alongside integration with bid-distribution methods
where those data are available (\S\ref{sec:forensic}).

\subsection{External validity and portability}
\label{sec:lim_external}

The empirical evidence comes from a single jurisdiction and a
single time window. The construct's portability rests on three
structural features explicit in BEC and identifiable in any
candidate jurisdiction: a commodity-heavy item composition that
supports within-product-code price comparisons; an electronic
platform with publicly auditable participant lists; and an
institutional asymmetry between winner and loser sides that makes
persistent zero wins informative about strategic non-competitive
participation. The architectural test in \S\ref{sec:forensic}
sharpens the portability claim: award-only enforcement
environments---those whose operational layer preserves
participants and winners but not bid-by-bid records---are the
construct's natural deployment environment, real-time-electronic
environments (preg\~ao-style auctions) are the regime in which
the screening signal is sharpest, and sealed-bid environments
without real-time interaction are the scope-restricted boundary
rather than the natural deployment one. Three settings the paper
does not test: the federal Brazilian procurement registry
(ComprasNet) and other state-level registries preserve the same
award-record envelope as BEC and are accessible through the same
legal regime; non-commodity settings (services, infrastructure,
public works) are within the framework's logical scope but the
within-item identifying variation that underwrites the headline
price gap is not equally available; and time periods outside the
sample window are untested. The $2021$ reform
(Lei~$14.133/2021$) consolidates preg\~ao-style auctions as the
institutional default---combined with the modal contrast in
\S\ref{sec:mech_modal}, the regulatory direction is favorable to
portability rather than adversarial.
